\documentclass[11pt]{article}

\usepackage[T1]{fontenc}
\usepackage[utf8]{inputenc}
\usepackage{graphicx}
\usepackage[margin=1in]{geometry}

\title{Agentic Software Observability:\
       Capturing, Transferring, and Re-executing Diagnostic Expertise}
\author{Yuvraj Sehgal\
        \small Department of Computer Science, Brock University\
        \small Supervisor: Dr.\ Naser Ezzati-Jivan}
\date{September 2026}

\begin{document}

\maketitle

\begin{abstract}
Software systems are increasingly diagnosed by AI agents rather than by human engineers. Diagnosing an issue across metrics, logs, distributed traces, and kernel events requires determining which signals isolate a root cause, which events warrant collection, and which thresholds matter for a specific system, and that determination is system-specific expertise that is rarely retained: it stays tacit knowledge held by individual engineers, or is discarded once an incident is resolved. Current approaches do not close this gap. Automated root-cause analysis localises faults from metrics, logs, and application traces, and recent LLM-based agents perform the same task with tool access, but in both the reusable asset is a model or a prompt; the reasoning that resolved one incident is not preserved in a form a later investigation can run, so an agent rebuilds each one from first principles, with little insight into how faults manifest below the application layer, a boundary at which the datasets these methods are evaluated on also stop. What is missing is expertise that survives an incident in executable form, evidence that it transfers to an unseen service or system, and a path by which a diagnosis determines what telemetry is collected next. This research develops \emph{agentic software observability}: how an autonomous agent acquires, retains, and applies diagnostic expertise derived from past incidents and historical root-cause analyses. An \emph{observability blueprint} makes the investigation itself durable, recording the condition it applies to, the telemetry needed to observe it, the analysis as runnable steps, the decision rule that produces a verdict, and the boundaries under which it must not be trusted. Blueprints accumulate into a curated library an agent selects from and executes unaided, admitted only once their signals are measured against every issue class under study rather than the one they target alone, and surrounded by the components this form of observability requires: evidence-based selection, a hierarchy narrowing from issue class to cause, and a feedback path from diagnosis into collection. Across roughly ten classes of operational issue and microservice, monolithic, and agentic architectures, comparative evaluation runs the same agent on the same incidents with and without the library against a strong deterministic baseline, on a dataset that aligns kernel events with metrics, logs, and distributed traces under labelled fault injection. It establishes whether a blueprint transfers across services and systems, whether the library abstains on uncovered problems, whether executable steps outperform prose, and whether evidence alone suffices to select the right blueprint, and characterises where captured expertise accelerates diagnosis, where a capable agent succeeds without it, and where it misleads, moving trace analysis from work that is repeated toward work that accumulates.
\end{abstract}

% ---------------------------------------------------------------------------
% SHORTER VARIANT, one paragraph, for a 300-350 word cap. Same moves: context,
% problem, what current methods do, the gap, "this research ...", then scope
% and evaluation. Swap by commenting out the version above.
%
% Software systems are increasingly diagnosed by AI agents rather than by human engineers. Diagnosing an issue across metrics, logs, distributed traces, and kernel events requires determining which signals isolate a root cause, which events warrant collection, and which thresholds matter for a specific system, and that expertise is rarely retained: it stays tacit knowledge held by individual engineers, or is discarded once an incident is resolved. Current approaches do not close this gap. Automated root-cause analysis localises faults from metrics, logs, and application traces, and recent LLM-based agents perform the same task with tool access, but in both the reusable asset is a model or a prompt; the reasoning that resolved one incident is not preserved in a form a later investigation can run, and the datasets these methods are evaluated on stop at the application boundary. What is missing is expertise that survives an incident in executable form, evidence that it transfers to an unseen system, and a path from diagnosis back into collection. This research develops \emph{agentic software observability}: how an autonomous agent acquires, retains, and applies diagnostic expertise derived from past incidents and historical root-cause analyses. An \emph{observability blueprint} makes the investigation itself durable, recording the condition it applies to, the telemetry needed, the analysis as runnable steps, the decision rule, and the boundaries under which it must not be trusted; blueprints accumulate into a curated library an agent selects from and executes unaided, admitted only once their signals are measured against every issue class under study. Across roughly ten classes of operational issue and microservice, monolithic, and agentic architectures, comparative evaluation runs the same agent with and without the library against a strong deterministic baseline, characterising where captured expertise accelerates diagnosis, where a capable agent succeeds without it, and where it misleads, moving trace analysis from work that is repeated toward work that accumulates.
% ---------------------------------------------------------------------------

\end{document}
